\section{Технический проект}

\subsection{Общая характеристика проектного решения}

Разрабатываемая система реализована по клиент-серверной архитектуре с разделением функциональности на серверную и клиентскую части. Конечный пользователь системы~--~сотрудник интернет-провайдера~--~работает с системой через современный браузер, не требующий установки никакого специализированного программного обеспечения; на стороне сервера размещены постоянное хранилище данных, бизнес-логика, проверки прав доступа и подсистема интеграции с маршрутизаторами MikroTik.

Серверная часть системы реализована на языке программирования PHP в составе фреймворка Laravel версии~12. Программный интерфейс прикладного уровня построен в стиле REST~\cite{rfc9110} с обменом данными в формате JSON по протоколу HTTP. Аутентификация запросов реализована средствами библиотеки Laravel~Sanctum: после успешного входа сотрудник получает токен Bearer, передаваемый в заголовке \texttt{Authorization} каждого последующего запроса.

Постоянное хранилище данных~--~реляционная система управления базами данных MySQL версии~8.0~\cite{mysqldoc}. Доступ к данным осуществляется через объектно-реляционное отображение Eloquent~ORM~\cite{romanov}, входящее в состав Laravel; запросы строятся через объекты моделей и построитель запросов, что исключает возможность внедрения SQL-кода~\cite{owasp} через значения пользовательского ввода.

Клиентская часть системы реализована как одностраничное веб-приложение (Single Page Application, SPA) на основе библиотеки Svelte версии~5 и метафреймворка SvelteKit версии~2~\cite{sveltedoc,sveltekitdoc}. Стилевое оформление построено поверх Tailwind~CSS версии~4~\cite{tailwind} с применением темы Cerberus библиотеки компонентов Skeleton~UI; цвета базовой палитры вынесены в CSS-переменные, что позволяет в дальнейшем расширить систему пользовательскими цветовыми темами без модификации компонентов. Печатные формы счетов на оплату формируются непосредственно на стороне клиента средствами библиотеки jsPDF.

В состав системы дополнительно интегрировано облачное файловое хранилище Cloudinary для размещения файловых вложений к обращениям технической поддержки и прикрепляемых сотрудниками документов-подтверждений ручных операций по лицевым счетам абонентов.

\emph{Замечание о границах работы.} Личный кабинет абонента, предоставляющий конечному потребителю услуг самостоятельный доступ к системе, в настоящей работе не реализуется и рассматривается как направление дальнейшего развития системы (см. заключение). Архитектура серверной части и схема базы данных при этом спроектированы с учетом возможности дальнейшего подключения такого личного кабинета и интеграции с произвольным платежным шлюзом для онлайн-пополнения баланса без переработки существующего ядра системы.

Общая схема организации системы с указанием всех ее ключевых составляющих и каналов взаимодействия между ними приведена на рисунке~\ref{fig:system_overview}.

\begin{figure}[ht]
\centering
\includegraphics[width=\linewidth,keepaspectratio]{diagrama_3_1}
\caption{Общая схема организации системы}
\label{fig:system_overview}
\end{figure}

При работе системы каждый запрос пользователя в браузере преобразуется в HTTP-запрос к серверной части по защищенному каналу HTTPS. На стороне сервера запрос принимается обратным прокси-сервером Nginx, который передает его на исполнение в PHP-FPM, где запрос обрабатывается фреймворком Laravel: маршрутизатор определяет соответствующий контроллер, промежуточные обработчики проверяют действующий токен и наличие необходимых прав, контроллер выполняет требуемую операцию через слой моделей и сервисов, обращаясь по необходимости к СУБД MySQL. Длительные операции (синхронизация очередей MikroTik, формирование счетов, рассылка уведомлений) не выполняются в рамках обработки HTTP-запроса; вместо этого контроллер ставит соответствующую задачу в очередь Laravel, а отдельный процесс-рабочий (worker) подхватывает задачу из очереди и выполняет ее асинхронно. Запуск фоновых задач по расписанию выполняется планировщиком Laravel, вызываемым ежеминутно через системный \texttt{cron}.

В режиме локальной разработки серверная часть запускается через встроенный сервер разработки Laravel (\texttt{php artisan serve}), клиентская часть~--~через сервер разработки Vite с поддержкой горячей замены модулей. В производственной конфигурации система развертывается в Docker-контейнерах под управлением платформы Coolify, обеспечивающей автоматическую сборку контейнеров из исходного кода в Git-репозитории, выпуск SSL-сертификатов через Let's~Encrypt и развертывание обновлений по принципу <<rolling update>>; клиентский код собирается командой \texttt{vite build} в каталог \texttt{build/} и отдается тем же веб-сервером Nginx, что и API-запросы.

\subsection{Выбор средств и технологий реализации}

\subsubsection{Состав применяемого технологического стека}

Для реализации системы выбран следующий набор инструментов и библиотек, перечисленных в порядке принадлежности к слоям приложения:

\emph{Серверная часть}:
\begin{itemize}
\item язык программирования PHP версии~8.2 или новее~\cite{phpdoc};
\item веб-фреймворк Laravel~12 (маршрутизация запросов, ORM, миграции, сеансы, защита от подделки запросов, очереди задач, планировщик)~\cite{laraveldoc};
\item библиотека Laravel~Sanctum~4 (аутентификация по токенам Bearer)~\cite{sanctum};
\item клиент RouterOS~API \texttt{evilfreelancer/routeros-api-php}~1.6 (взаимодействие с маршрутизаторами MikroTik)~\cite{routerosphp};
\item клиент Cloudinary~--~\texttt{cloudinary-labs/cloudinary-laravel}~3.0 (загрузка файлов в облачное хранилище);
\item тестовый фреймворк PestPHP~4.1 (модульное и функциональное тестирование).
\end{itemize}

\emph{Клиентская часть}:
\begin{itemize}
\item библиотека Svelte~5 (реактивные компоненты пользовательского интерфейса)~\cite{sveltedoc};
\item метафреймворк SvelteKit~2 (маршрутизация на стороне клиента, сборка)~\cite{sveltekitdoc};
\item CSS-фреймворк Tailwind~CSS~4 в сочетании с библиотекой компонентов Skeleton~UI~4~\cite{tailwind};
\item библиотеки jsPDF~3.0 и jspdf-autotable~5.0 (формирование печатных форм);
\item библиотека Lucide~Icons (\texttt{@lucide/svelte}~0.546)~--~векторные пиктограммы;
\item библиотека svelte-sonner~1.0 (всплывающие уведомления);
\item библиотека recharts~3.3 (построение графиков на дашборде);
\item сборщик Vite~7.1 и тестовый фреймворк Vitest~4.1.
\end{itemize}

\emph{СУБД и инфраструктура}: MySQL~8.0~\cite{mysqldoc} в качестве постоянного хранилища; Docker в качестве среды выполнения контейнеров; Nginx в качестве обратного прокси-сервера и сервера статических файлов; Coolify в качестве платформы непрерывного развертывания.

Выбор технологий основан на четырех соображениях: \emph{зрелость} (все перечисленные позиции применяются в коммерческих проектах годами); \emph{лицензионная чистота} (открытый исходный код по разрешительным лицензиям MIT, BSD, ISC, Apache~2.0, без платежей в иностранной валюте); \emph{доступность документации и сообщества}; \emph{совместимость со стандартными протоколами} (HTTP/JSON, MySQL, RouterOS~API), что позволяет заменить любую часть на аналог.

\subsubsection{Серверный язык программирования и фреймворк Laravel}

Язык PHP выбран как исторически разработанный для веб-программирования, обладающий наиболее обширной поддержкой шаблонизации, обработки HTTP-запросов и взаимодействия с базами данных~\cite{phpdoc}. Производительность PHP в режиме PHP-FPM достаточна для оператора малого/среднего масштаба: основной вклад во время отклика вносят запросы к базе и обращения к маршрутизатору, доля интерпретации PHP в общем времени отклика не превышает 10--15\,\%.

Веб-фреймворк Laravel~12~\cite{laraveldoc} к своей последней версии аккумулировал практически весь спектр функциональности типового веб-приложения: объектно-реляционное отображение Eloquent~ORM с автоматическими миграциями~\cite{romanov}; штатную систему аутентификации с хешированием Bcrypt; защиту от CSRF/XSS/кликджекинга в промежуточных обработчиках; валидацию входных данных через Form Requests; подсистему очередей с несколькими драйверами и встроенным планировщиком; обширную экосистему пакетов через Composer.

\subsubsection{Аутентификация через Laravel~Sanctum}

Библиотека Laravel~Sanctum~\cite{sanctum} реализует токенную аутентификацию для одностраничных приложений и мобильных клиентов. Альтернатива Passport (полнофункциональный OAuth~2.0) для разрабатываемой системы избыточна. При успешном входе сервер создает запись токена в \texttt{personal\_access\_tokens} (хеш SHA-256), клиент сохраняет открытый текст токена в \texttt{localStorage} и прикладывает в заголовке \texttt{Authorization: Bearer~...} ко всем запросам. Параллельно Sanctum реализует CSRF-токены: при изменяющих запросах клиент передает значение токена в заголовке \texttt{X-CSRFToken}.

\subsubsection{Клиентская библиотека Svelte и метафреймворк SvelteKit}

Svelte~\cite{sveltedoc}~--~компилируемая библиотека реактивных компонентов. Принципиальное отличие от React и Vue в том, что Svelte~--~компилятор: компоненты на этапе сборки превращаются в JavaScript-код прямых императивных операций с DOM. Это дает меньший размер итогового бандла (в 2--3 раза меньше, чем React, за счет отсутствия runtime-библиотеки) и более высокую производительность при множественных обновлениях состояния. С версии~5 применяется модель <<рун>> с явным объявлением реактивных значений через \texttt{\$state()}, \texttt{\$derived()}, \texttt{\$effect()}.

Метафреймворк SvelteKit~\cite{sveltekitdoc} добавляет файловую маршрутизацию (каждый файл в \texttt{src/routes/} соответствует маршруту), серверный рендеринг (опционально), оптимизацию загрузки данных через функции \texttt{load()}, интеграцию со сборщиком Vite. Для системы выбран режим без серверного рендеринга (приложение работает за аутентификацией и не индексируется); собранный код~--~набор статических файлов, не требующий процесса Node.js.

\subsubsection{СУБД MySQL и библиотека RouterOS~API}

В качестве хранилища выбрана СУБД MySQL~8.0~\cite{mysqldoc}~--~одна из самых распространенных свободных реляционных СУБД. Из специфических возможностей применены: проверочные ограничения CHECK на уровне таблицы; генерируемые столбцы (\texttt{full\_name} в таблицах абонентов и сотрудников); тип JSON с выборкой по содержимому (для \texttt{metadata} транзакций, \texttt{configuration\_snapshot} счетов, \texttt{old\_values}/\texttt{new\_values} аудита); кодировка \texttt{utf8mb4} для полной поддержки Unicode. Управление схемой~--~через миграции Laravel; ручное редактирование таблиц запрещено.

Для взаимодействия с маршрутизаторами по протоколу RouterOS~API применяется библиотека \texttt{evilfreelancer/routeros-api-php}~1.6~\cite{routerosphp} (лицензия MIT, доступна через Composer). Библиотека реализует бинарную кодировку сообщений, аутентификацию, поддержку TLS, асинхронных команд и потоковых подписок. Альтернатива \texttt{routeros4me/routeros-api-php} реже обновляется (последний релиз~--~2021~г.), что и определило выбор.

\subsection{Архитектурное построение системы}

\subsubsection{Уровневая модель архитектуры}

Программный код разрабатываемой системы организован по принципу разделения ответственности на уровни (англ. layers)~\cite{fowler_enterprise,orlov}. Каждый уровень имеет четко очерченную область ответственности и обращается только к смежному с ним нижестоящему уровню; такая организация уменьшает связность компонентов и упрощает локализацию ошибок при разработке и сопровождении. Уровневая схема архитектуры разрабатываемой системы приведена на рисунке~\ref{fig:layered_arch}.

\begin{figure}[H]
\centering
\includegraphics[width=0.7\linewidth,height=0.85\textheight,keepaspectratio]{diagrama_3_2}
\caption{Уровневая схема архитектуры системы}
\label{fig:layered_arch}
\end{figure}

Из шести представленных уровней первые два работают непосредственно в браузере конечного пользователя, четыре следующих~--~на сервере. Граница между клиентом и сервером проходит по HTTP-вызову. На клиенте уровень обмена упаковывает запрос в HTTP-сообщение и отправляет его на сервер, где уровень программного интерфейса принимает запрос и разбирает его. Между уровнем бизнес-логики и уровнем хранения встроено объектно-реляционное отображение Eloquent~ORM, поэтому ни одна строка прикладного кода не собирает SQL-выражения вручную: запросы строятся через цепочки методов построителя запросов либо через объектно-ориентированные методы моделей.

\subsubsection{Компоненты клиентской и серверной частей}

Код клиентской части (\texttt{src/} проекта SvelteKit) разделен на три группы: \emph{компоненты} (\texttt{src/lib/components/})~--~многоразовые элементы интерфейса (базовые \texttt{Button}, \texttt{Input}, \texttt{Modal} и тематические \texttt{Menu}, \texttt{Chart}, \texttt{KpiCards}, \texttt{ModalVerFactura}, \texttt{FirewallTable}); \emph{страницы} (\texttt{src/routes/})~--~сборочные единицы маршрутов SvelteKit, перечисленные в таблице~\ref{tbl:routes}; \emph{хранилища состояния} (\texttt{src/lib/stores/})~--~глобальное состояние на основе рун Svelte~5: \texttt{auth.svelte.ts} (токен, имя, роль, разрешения, метод \texttt{can(module,action)}), \texttt{app.svelte.ts} (UI-состояние), \texttt{usuarios.svelte.ts} (кэш сотрудников), \texttt{mikrotik-firewall.svelte.ts} (локальный черновик правил брандмауэра).

\begin{xltabular}{\textwidth}{|p{4.2cm}|X|}
\caption{Маршруты клиентской части системы\label{tbl:routes}}\\ \hline
\centrow Маршрут & \centrow Назначение страницы \\ \hline
\endfirsthead
\continuecaption{Продолжение таблицы \ref{tbl:routes}}
\centrow Маршрут & \centrow Назначение страницы \\ \hline
\finishhead
\route{/login} & Страница входа в панель администрирования (форма ввода адреса электронной почты и пароля) \\ \hline
\route{/} & Главная страница (дашборд): ключевые показатели, графики, перечень должников \\ \hline
\route{/clientes} & Реестр абонентов с боковой панелью списка и центральной карточкой выбранного абонента \\ \hline
\route{/planes} & Каталог тарифных планов с модальным окном создания и редактирования \\ \hline
\route{/facturas} & Реестр счетов на оплату с фильтрами по периоду, статусу и абоненту \\ \hline
\route{/proveedores} & Реестр внешних интернет-провайдеров (поставщиков канала) \\ \hline
\route{/proveedores/[id]/conexiones} & Реестр подключений к конкретному поставщику \\ \hline
\route{/usuarios} & Реестр сотрудников панели администрирования \\ \hline
\route{/usuarios/roles} & Управление ролями и матрица разрешений \\ \hline
\route{/configuraciones} & Хаб системных настроек \\ \hline
\route{/configuraciones/facturacion} & Настройки фискальных реквизитов выпуска счетов \\ \hline
\route{/configuraciones/importacion} & Импорт абонентов и тарифных планов из CSV \\ \hline
\route{/configuraciones/workers} & Управление расписанием автоматических задач \\ \hline
\route{/mikrotik} & Обзорная страница раздела управления оборудованием MikroTik \\ \hline
\route{/mikrotik/dispositivos} & Реестр маршрутизаторов в системе \\ \hline
\route{/mikrotik/firewall} & Управление правилами брандмауэра выбранного маршрутизатора \\ \hline
\route{/mikrotik/colas} & Управление простыми очередями и их мониторинг \\ \hline
\route{/mikrotik/monitoreo} & Мониторинг подключенных беспроводных устройств \\ \hline
\route{/mikrotik/sincronizacion} & Управление расписанием синхронизации очередей и журнал результатов \\ \hline
\route{/perfil} & Просмотр и редактирование собственного профиля пользователя
\end{xltabular}

Код серверной части (\texttt{app/} проекта Laravel) разделен на подкаталоги по категориям классов:
\begin{itemize}
\item \texttt{app/Models/}~--~описания таблиц базы данных как классов-наследников \texttt{Illuminate\textbackslash Database\textbackslash Eloquent\textbackslash Model}; в системе насчитывается 26 моделей, охватывающих все ключевые сущности предметной области;
\item \texttt{app/Http/Controllers/}~--~контроллеры, обрабатывающие HTTP-запросы; основная часть размещена в \texttt{app/Http/Controllers/Admin/} и обслуживает эндпоинты панели администрирования;
\item \texttt{app/Http/Requests/}~--~классы валидации входных данных (Form Requests); каждый запрос с телом имеет соответствующий класс с описанием правил валидации;
\item \texttt{app/Http/Resources/}~--~сериализаторы (API Resources), преобразующие модели Eloquent в JSON-структуры ответа сервера; применяются для контроля состава полей в ответе и сокрытия внутренних полей (например, паролей);
\item \texttt{app/Http/Middleware/}~--~промежуточные обработчики запросов; среди них \texttt{EnsureEmployeeSuperAdmin} (проверка роли суперадминистратора), \texttt{CheckPermission} (проверка наличия указанного разрешения у текущего пользователя), \texttt{AuditableMiddleware} (фиксация запроса в журнале аудита);
\item \texttt{app/Services/}~--~сервисы, инкапсулирующие бизнес-логику, не привязанную к конкретному запросу: \texttt{MikroTikService} (низкоуровневые операции с маршрутизаторами), \texttt{MikroTikQueueSyncService} (алгоритм синхронизации очередей), \texttt{AutoBillingService} (формирование счетов и автоматическая оплата), \texttt{ClientSuspensionService} (приостановка и восстановление обслуживания), \texttt{AutomationSettingsService} (чтение настроек планировщика из БД);
\item \texttt{app/Jobs/}~--~фоновые задачи, выполняемые рабочим процессом очереди: \texttt{GenerateMonthlyInvoices}, \texttt{ProcessClientSuspension}, \texttt{ProcessAutoReactivation}, \texttt{SyncMikroTikQueues}, \texttt{ApplyFirewallChanges};
\item \texttt{app/Console/Commands/}~--~собственные консольные команды Artisan: \texttt{TestMikroTikConnection}, \texttt{CheckAndReactivate}, \texttt{ProcessAutoBilling}.
\end{itemize}

Маршрутизация запросов определена в файлах \texttt{routes/api.php} (программный интерфейс прикладного уровня), \texttt{routes/console.php} (расписание автоматических задач), \texttt{routes/web.php} (одна-единственная страница приветствия для невалидных обращений). Маршруты в \texttt{routes/api.php} автоматически получают префикс \texttt{/api} и сгруппированы по подсистемам с применением вложенных групп middleware: внешняя группа применяет middleware \texttt{auth:sanctum} ко всем маршрутам, кроме входа в систему; вложенные группы применяют middleware проверки конкретных разрешений к маршрутам, относящимся к одной подсистеме.

\subsubsection{Компоненты уровня доступа к данным}

Уровень доступа к данным разрабатываемой системы реализован средствами объектно-реляционного отображения Eloquent~ORM~\cite{romanov}, входящего в состав ядра Laravel. Каждая таблица базы данных описывается классом-наследником \texttt{Illuminate\textbackslash Database\textbackslash Eloquent\textbackslash Model}; поля таблицы являются атрибутами объекта модели, доступными через стандартные операторы PHP. Связи между таблицами (один к одному, один ко многим, многие ко многим, полиморфные связи) задаются методами класса модели, возвращающими специализированные объекты-связи (\texttt{HasOne}, \texttt{HasMany}, \texttt{BelongsTo}, \texttt{BelongsToMany}). Обращение к связанным записям выполняется лениво (по запросу через свойство-связь) либо энергично (через метод \texttt{with()} построителя запросов с автоматическим объединением запросов).

Ограничения целостности данных задаются как на уровне миграции (внешние ключи с указанием поведения при удалении связанной записи, уникальные индексы, проверочные ограничения), так и на уровне модели (правила валидации в Form Requests, события модели \texttt{creating}, \texttt{updating}, \texttt{deleting} с дополнительными проверками). Атомарные операции, затрагивающие несколько связанных таблиц одновременно, оборачиваются в одну транзакцию базы данных через статический метод \texttt{DB::transaction()} либо через построитель запросов \texttt{transaction.atomic}.

Часто исполняемые запросы оптимизируются методами \texttt{select\_related()} (для связей по внешнему ключу) и \texttt{prefetch\_related()} (для обратных отношений). Эти методы изменяют стратегию выборки: вместо отдельного запроса за связанной записью при каждом обращении к свойству-связи выполняется одно объединение всей коллекции записей с их связанными, что значительно сокращает количество SQL-запросов и общее время выполнения операции.

Для автоматического протоколирования изменений данных применяется trait \texttt{Auditable}, подключаемый к моделям, изменения в которых подлежат журнализации. При срабатывании событий модели \texttt{saving}, \texttt{updating}, \texttt{deleting} trait автоматически записывает в таблицу \texttt{audits} имя измененной таблицы, идентификатор измененной записи, тип операции, старое и новое значения измененных полей в формате JSON, идентификатор инициатора изменения (по данным текущего сеанса аутентификации), IP-адрес источника запроса.

\subsubsection{Диаграмма последовательности обработки запроса}

Для иллюстрации взаимодействия компонентов системы при обработке типового запроса пользователя рассмотрим сценарий <<Сотрудник приостанавливает обслуживание абонента из карточки абонента>>. Последовательность шагов обработки приведена на рисунке~\ref{fig:sequence_suspend}.

\begin{figure}[H]
\centering
\includegraphics[width=\linewidth,height=0.85\textheight,keepaspectratio]{diagrama_3_3}
\caption{Последовательность обработки запроса <<Приостановка обслуживания абонента>>}
\label{fig:sequence_suspend}
\end{figure}

Шаги обработки запроса:
\begin{enumerate}
\item Пользователь в браузере нажимает кнопку <<Приостановить обслуживание>> на карточке абонента в панели администрирования. Браузер передает событие нажатия компоненту Svelte, обрабатывающему карточку.
\item Компонент собирает параметры операции, формирует HTTP-запрос \texttt{POST /api/admin/clientes/\{id\}/suspend} с прикреплением токена Bearer в заголовке \texttt{Authorization} и значения CSRF-токена в заголовке \texttt{X-CSRFToken}, отправляет запрос на сервер.
\item Сервер принимает запрос; промежуточные обработчики Laravel последовательно выполняют проверку токена аутентификации (через Sanctum) и проверку наличия у текущего пользователя разрешения \texttt{clientes.editar}. При отсутствии токена или разрешения дальнейшая обработка прекращается с возвратом ответа 401 или 403 соответственно.
\item Контроллер \texttt{ClientController} получает идентификатор абонента из параметров маршрута, обращается к модели \texttt{Client} с методом \texttt{findOrFail(\$id)} для загрузки записи; модель формирует SQL-запрос и отправляет его в СУБД через драйвер PDO.
\item Контроллер вызывает метод \texttt{suspend()} сервиса \texttt{ClientSuspensionService}, который открывает транзакцию, изменяет статус абонента на \texttt{suspended}, сохраняет изменения.
\item Сервис создает запись в таблице \texttt{client\_events} с типом события \texttt{suspended}, фиксирующую состоявшееся изменение.
\item Сервис диспатчит задачу \texttt{SyncMikroTikQueues} в очередь Laravel с приоритетом \texttt{high}; контроллер возвращает клиенту ответ 200~OK с обновленной записью абонента. Транзакция базы данных коммитится.
\item Рабочий процесс (worker) подхватывает задачу из очереди и устанавливает соединение с маршрутизатором MikroTik абонента через сервис \texttt{MikroTikService}. Сначала задача выполняет команду удаления простой очереди абонента (\texttt{/queue/simple/remove}).
\item Затем задача выполняет команду добавления IP-адреса абонента в адресный список \texttt{suspended} (\texttt{/ip/firewall/address-list/add}).
\item Клиент сотрудника получает ответ 200~OK и обновляет карточку абонента с новым статусом. Изменение фиксируется в журнале аудита и доступно остальным сотрудникам при следующем обращении к карточке абонента.
\end{enumerate}

\subsubsection{Форматы обмена данными}

Все тела HTTP-запросов и ответов в программном интерфейсе системы передаются в формате JSON. Поля документов кодируются в кодировке UTF-8. Календарные значения (даты, отметки времени) передаются строками в формате ISO~8601 (например, \texttt{<<2026-05-16T18:30:00Z>>}); сервер при необходимости разбирает строку и сохраняет в столбец типа \texttt{DATETIME(6)}. Идентификаторы записей передаются целыми числами; денежные значения~--~строками с двумя знаками после десятичной точки (для предотвращения потерь точности при сериализации через \texttt{float}).

Списочные ресурсы возвращаются в постраничной обертке Laravel API~Resources: ответ представляет собой объект с полями \texttt{data} (массив записей текущей страницы), \texttt{links} (ссылки на первую, предыдущую, текущую, следующую и последнюю страницы), \texttt{meta} (метаинформация о пагинации: текущая страница, общее число страниц, общее число записей, размер страницы, диапазон номеров записей). По умолчанию страница содержит 25 записей; размер страницы регулируется параметром запроса \texttt{per\_page}.

При возникновении ошибок сервер возвращает HTTP-ответ с соответствующим кодом состояния и телом, описывающим ошибку:
\begin{itemize}
\item код 400 (Bad Request)~--~некорректный формат запроса (невалидный JSON, отсутствие обязательных параметров маршрута);
\item код 401 (Unauthorized)~--~отсутствие или невалидность токена аутентификации;
\item код 403 (Forbidden)~--~недостаточные права на выполнение операции;
\item код 404 (Not Found)~--~запрашиваемая запись или маршрут не существуют;
\item код 419 (Page Expired)~--~отсутствие или несовпадение CSRF-токена;
\item код 422 (Unprocessable Entity)~--~ошибки валидации полей; в теле ответа возвращается объект \texttt{errors} с массивами сообщений по каждому невалидному полю;
\item код 429 (Too Many Requests)~--~превышение лимита частоты обращений (rate limiting);
\item код 500 (Internal Server Error)~--~непредвиденная ошибка серверной логики; в производственной конфигурации детализированное сообщение об ошибке не возвращается клиенту, а заносится в системный журнал.
\end{itemize}

\subsubsection{Аутентификация и авторизация}

Аутентификация запросов работает через библиотеку Laravel~Sanctum~\cite{sanctum}, как подробно описано в подразделе~3.2.3. Параллельно с проверкой действующего токена сервер на каждом изменяющем запросе проверяет наличие у пользователя необходимого разрешения через middleware \texttt{permission}. Разрешение указывается в строке маршрута в формате \texttt{permission:модуль.действие}; например, маршрут создания нового счета определен как:

\begin{verbatim}
Route::post('/admin/invoices', [AdminInvoiceController::class, 'store'])
     ->middleware('permission:facturas.crear');
\end{verbatim}

Промежуточный обработчик \texttt{CheckPermission} извлекает текущего сотрудника из контекста запроса, обращается к методу \texttt{hasPermission(\$slug)} модели сотрудника, который, в свою очередь, проверяет наличие разрешения с указанным \texttt{slug} среди разрешений, привязанных к роли сотрудника через связующую таблицу \texttt{role\_permission}. Для сотрудников с ролью \texttt{super\_admin} проверка всегда возвращает положительный результат, что устраняет необходимость явного назначения всех разрешений суперадминистраторской роли.

Дополнительный уровень защиты обеспечивается политиками доступа (Policies) Laravel: для каждой модели, доступной через программный интерфейс системы, определен класс политики, проверяющий конкретные правила доступа на уровне записи. Например, политика \texttt{ClientPolicy} реализует метод \texttt{view(\$user, \$client)}, возвращающий \texttt{true}, если текущий пользователь либо является сотрудником с правом \texttt{clientes.ver}, либо является самим абонентом~--~владельцем учетной записи. Такая двухуровневая схема (общее разрешение на действие + проверка прав на конкретную запись) исключает возможность доступа сотрудников к записям, не относящимся к их зоне ответственности.

\subsubsection{Описание программного интерфейса}

Программный интерфейс системы размещен под префиксом \texttt{/api/} и сгруппирован по подсистемам. Общее число эндпоинтов~--~около 80; они охватывают аутентификацию (\texttt{/api/employee/login}, выход, получение CSRF-токена), CRUD по основным сущностям (абоненты, тарифные планы, подписки, счета, маршрутизаторы, поставщики канала, сотрудники, роли, разрешения), специализированные операции (изменение статуса абонента, автоматическое формирование счетов, синхронизация очередей, применение и откат правил брандмауэра), подсистему обращений (создание, ведение переписки сотрудниками, назначение исполнителя, изменение статуса), системные настройки и автоматизации, импорт данных и панель аналитики. Полная детализация эндпоинтов с примерами тел запросов и ответов вынесена в приложение.

Каждый эндпоинт защищен последовательностью промежуточных обработчиков: \texttt{auth:sanctum} (проверка действующего токена), \texttt{permission:модуль.действие} (проверка наличия разрешения у текущего пользователя), \texttt{throttle} (ограничение частоты обращений для эндпоинтов, потенциально подверженных перебору). Эндпоинты, требующие роли \texttt{super\_admin} (управление маршрутизаторами, ролями, системными настройками), защищены отдельным обработчиком \texttt{EnsureEmployeeSuperAdmin}.

\subsection{Проект базы данных}

\subsubsection{ER-модель данных}

Концептуальная модель данных системы построена по методике <<сущность~--~связь>> (Entity-Relationship)~\cite{date,connolly}. Выделены тридцать сущностей, сгруппированных по шести функциональным подсистемам:
\begin{enumerate}
\item \emph{Управление абонентами}: \texttt{Client} (центральная сущность~--~абонент), \texttt{ClientPlan} (подписка абонента на тарифный план), \texttt{Plan} (тарифный план), \texttt{PlanFeature} (особенность тарифного плана для отображения в каталоге).
\item \emph{Финансы}: \texttt{Invoice} (счет на оплату), \texttt{Wallet} (лицевой счет абонента), \texttt{Transaction} (операция по лицевому счету).
\item \emph{Подсистема поддержки}: \texttt{Ticket} (обращение в техническую поддержку), \texttt{Message} (сообщение во внутренней переписке сотрудников по обращению), \texttt{Attachment} (файловое вложение к сообщению), \texttt{ClientEvent} (событие, относящееся к абоненту).
\item \emph{Доступ}: \texttt{User} (унаследованная техническая модель Laravel), \texttt{Employee} (сотрудник), \texttt{Role} (роль), \texttt{Permission} (разрешение), \texttt{role\_permission} (связующая таблица многие-ко-многим).
\item \emph{Сетевая инфраструктура}: \texttt{MikrotikRouter} (маршрутизатор в реестре), \texttt{FirewallFilterRule} (правило брандмауэра filter), \texttt{FirewallNatRule} (правило NAT), \texttt{FirewallApplyLog} (журнал применений правил), \texttt{InternetServiceProvider} (внешний поставщик канала), \texttt{IspConnection} (подключение к поставщику).
\item \emph{Конфигурация и инфраструктурные сущности}: \texttt{Setting} (системная настройка), \texttt{AutomationSetting} (настройка автоматической задачи), \texttt{ImportHistory} (журнал импортов), \texttt{Audit} (запись журнала аудита).
\end{enumerate}

Ключевые связи между сущностями: \texttt{Client} 1:N \texttt{ClientPlan}, \texttt{Plan} 1:N \texttt{ClientPlan} (треугольник <<абонент~--~подписка~--~тариф>>); \texttt{Client} 1:1 \texttt{Wallet} 1:N \texttt{Transaction}; \texttt{Client} 1:N \texttt{Invoice} (счет может ссылаться на конкретную подписку через необязательный внешний ключ \texttt{client\_plan\_id}); \texttt{Client} 1:N \texttt{Ticket} 1:N \texttt{Message} 1:N \texttt{Attachment}; \texttt{Employee} N:1 \texttt{Role}, \texttt{Role} N:M \texttt{Permission} (через таблицу \texttt{role\_permission}); \texttt{MikrotikRouter} 1:N \texttt{FirewallFilterRule}, \texttt{FirewallNatRule}, \texttt{FirewallApplyLog}; \texttt{InternetServiceProvider} 1:N \texttt{IspConnection}.

Обобщенная ER-диаграмма основных сущностей системы приведена на рисунке~\ref{fig:er_model}. В целях компактности на диаграмме отображены только ключевые сущности из подсистем 1--3; модель доступа, сетевой инфраструктуры и конфигурации рассмотрены отдельно в текстовом изложении.

\begin{figure}[ht]
\centering
\includegraphics[width=\linewidth,height=0.85\textheight,keepaspectratio]{diagrama_3_4}
\caption{ER-модель основных сущностей системы (подсистемы абонентов, финансов, поддержки)}
\label{fig:er_model}
\end{figure}

\subsubsection{Нормализация данных}

Схема базы данных приведена к третьей нормальной форме~\cite{date}.

Первая нормальная форма достигается автоматически: в таблицах нет атрибутов-списков и повторяющихся групп значений. Цвета, скорости и прочие параметры разложены по отдельным полям; JSON-структуры применяются только в полях, явно предназначенных для хранения нетипизированного содержимого (\texttt{metadata} в транзакциях, \texttt{configuration\_snapshot} в счетах, \texttt{old\_values}/\texttt{new\_values} в журнале аудита).

Вторая нормальная форма достигается за счет применения суррогатных первичных ключей: все таблицы построены на автоинкрементном поле \texttt{id}, поэтому частичной зависимости атрибутов от части составного ключа в принципе быть не может. Случаи, где напрашивался бы составной ключ (пара \texttt{client\_id, plan\_id} в подписках; пара \texttt{predecessor\_id, successor\_id} в зависимостях; пара \texttt{role\_id, permission\_id} в связке ролей и разрешений), реализованы как ограничения уникальности на индексах при сохранении суррогатного ключа.

Третья нормальная форма потребовала отдельной проверки. В таблице абонентов хранится производное поле \texttt{full\_name}~--~объединение фамилии, имени и отчества; формально оно зависит транзитивно от других полей таблицы. Это нарушение устранено за счет применения механизма генерируемых столбцов MySQL: значение поля \texttt{full\_name} объявлено как \texttt{GENERATED ALWAYS AS (CONCAT\_WS(' ', last\_name, first\_name, middle\_name)) VIRTUAL}, что фактически делает поле производным от других на уровне СУБД, а не отдельно хранимым. Аналогичное решение применено для поля \texttt{full\_name} таблицы сотрудников. В остальных таблицах транзитивные зависимости устранены: в таблице задач хранятся только идентификаторы абонента и тарифа, а не их данные; в таблице участников не дублируется имя пользователя или название проекта.

\emph{Особый случай~--~снимки в журнале счетов и журнале аудита}: поле \texttt{configuration\_snapshot} таблицы счетов хранит JSON-копию фискальных реквизитов организации на момент выпуска счета, а поля \texttt{old\_values}, \texttt{new\_values} таблицы аудита~--~JSON-снимки значений измененных полей. Такие денормализованные представления не нарушают третью нормальную форму, поскольку семантика этих полей~--~именно <<заморозить>> данные в виде, актуальном на момент события, а не отразить текущее состояние ссылочной записи. Изменения в источнике после момента события не должны отражаться на сохраненной копии~--~это требование журнала.

\subsubsection{Описание структуры базы данных}

Полная физическая схема базы данных системы со всеми прикладными и служебными таблицами, типами полей и внешними ключами приведена на рисунке~\ref{fig:db_schema}.

\begin{landscape}
\thispagestyle{plain}
\begin{figure}[!h]
\centering
\includegraphics[width=27cm,height=15.5cm,keepaspectratio]{diagrama_db.png}
\caption{Физическая схема базы данных системы}
\label{fig:db_schema}
\end{figure}
\end{landscape}

Физическая схема базы данных системы состоит из прикладных таблиц и служебных таблиц фреймворка Laravel. К прикладным относятся:
\texttt{clients}, \texttt{plans}, \texttt{plan\_features}, \texttt{clients\_plans},
\texttt{invoices}, \texttt{wallets}, \texttt{transactions},
\texttt{tickets}, \texttt{messages}, \texttt{attachments}, \texttt{client\_events},
\texttt{employees}, \texttt{roles}, \texttt{permissions}, \texttt{role\_permission},
\texttt{mikrotik\_routers}, \texttt{firewall\_filter\_rules}, \texttt{firewall\_nat\_rules}, \texttt{firewall\_apply\_logs},
\texttt{internet\_service\_providers}, \texttt{isp\_connections},
\texttt{system\_settings}, \texttt{automation\_settings}, \texttt{import\_history}, \texttt{audits}.
Остальные таблицы создаются встроенными подсистемами Laravel
(\texttt{users}, \texttt{personal\_access\_tokens}, \texttt{jobs}, \texttt{failed\_jobs},
\texttt{cache}, \texttt{sessions}, \texttt{migrations}, \texttt{password\_reset\_tokens}, \texttt{notifications}).

Все прикладные модели используют тип первичного ключа \texttt{BIGINT UNSIGNED AUTO\_INCREMENT}; текстовые описания хранятся как \texttt{LONGTEXT}, отметки времени~--~как \texttt{DATETIME(6)}, денежные значения~--~как \texttt{DECIMAL(12,2)}. Связи между прикладными таблицами реализованы внешними ключами с поведением при удалении: \texttt{CASCADE} для строгих связей (подписка~--~абонент), \texttt{SET~NULL} для слабых (сотрудник-исполнитель обращения), \texttt{RESTRICT} для ссылок на создателя финансовых документов.

Обобщенные сведения о назначении прикладных таблиц приведены в таблице~\ref{tbl:db_summary}; ниже подробно рассматривается структура нескольких ключевых таблиц.

\begin{xltabular}{\textwidth}{|p{4.2cm}|X|}
\caption{Назначение прикладных таблиц базы данных\label{tbl:db_summary}}\\ \hline
\centrow Таблица & \centrow Назначение \\ \hline
\endfirsthead
\continuecaption{Продолжение таблицы \ref{tbl:db_summary}}
\centrow Таблица & \centrow Назначение \\ \hline
\finishhead
\texttt{clients} & Реестр абонентов: ФИО, документ, контакты, адрес установки, координаты, IP, статус обслуживания, дата заключения договора, баланс, хеш пароля абонента \\ \hline
\texttt{plans} & Каталог тарифных планов: наименование, скорости загрузки/отдачи, ежемесячная стоимость, расчетный период, имя очереди MikroTik, признак активности \\ \hline
\texttt{plan\_features} & Особенности тарифа для отображения в публичном каталоге \\ \hline
\texttt{clients\_plans} & Подписки абонента на тарифный план: даты, текущая цена, идентификатор очереди MikroTik, IP-адрес, статус \\ \hline
\texttt{invoices} & Счета на оплату: уникальный фискальный номер, суммы с налогом и без, статус, метод оплаты, снимок фискальных реквизитов \\ \hline
\texttt{wallets} & Лицевые счета абонентов: баланс, валюта, статус \\ \hline
\texttt{transactions} & Операции по лицевому счету: пополнение, оплата, возврат; прикрепленный документ-подтверждение в Cloudinary \\ \hline
\texttt{tickets} & Обращения в техническую поддержку: автор, исполнитель, статус, оценка \\ \hline
\texttt{messages} & Сообщения переписки сотрудников по обращению с поддержкой системных событий \\ \hline
\texttt{attachments} & Файловые вложения сообщений, размещенные в Cloudinary \\ \hline
\texttt{client\_events} & События обслуживания абонента (приостановка, восстановление, смена тарифа) для журнализации и аудита \\ \hline
\texttt{employees} & Учетные записи сотрудников панели администрирования с привязкой к роли \\ \hline
\texttt{roles}, \texttt{permissions}, \texttt{role\_permission} & Ролевая модель разграничения доступа: роли, разрешения формата \texttt{<<модуль.действие>>}, связующая таблица многие-ко-многим \\ \hline
\texttt{mikrotik\_routers} & Реестр маршрутизаторов MikroTik с зашифрованными учетными данными \\ \hline
\texttt{firewall\_filter\_rules} & Правила фильтрации трафика брандмауэра MikroTik \\ \hline
\texttt{firewall\_nat\_rules} & Правила трансляции сетевых адресов брандмауэра MikroTik \\ \hline
\texttt{firewall\_apply\_logs} & Журнал применений изменений правил со снимками <<до>> и <<после>> для возможности отката \\ \hline
\texttt{internet\_service\_providers}, \texttt{isp\_connections} & Реестр внешних поставщиков канала связи и параметры подключений \\ \hline
\texttt{system\_settings} & Типизированные системные настройки модулей (биллинг, фискальные реквизиты, MikroTik) \\ \hline
\texttt{automation\_settings} & Расписание автоматических задач планировщика \\ \hline
\texttt{import\_history} & Журнал массовых импортов с возможностью отката \\ \hline
\texttt{audits} & Сквозной журнал аудита всех изменений данных
\end{xltabular}

Ниже приведены подробные структуры таблиц, играющих центральную роль в системе. Структура реестра абонентов \texttt{clients} приведена в таблице~\ref{tbl:db_clients}.

\begin{xltabular}{\textwidth}{|p{3.0cm}|p{3.0cm}|p{1.6cm}|X|}
\caption{Структура таблицы \texttt{clients}\label{tbl:db_clients}}\\ \hline
\centrow Поле & \centrow Тип & \centrow Обязат. & \centrow Назначение \\ \hline
\endfirsthead
\continuecaption{Продолжение таблицы \ref{tbl:db_clients}}
\centrow Поле & \centrow Тип & \centrow Обязат. & \centrow Назначение \\ \hline
\finishhead
id & BIGINT & Да & Первичный ключ \\ \hline
\dbfield{full\_name} &VARCHAR(452) & Да & Полное имя (генерируемое поле) \\ \hline
\dbfield{document\_id} &VARCHAR(20) & Да & Документ, удостоверяющий личность (уникален) \\ \hline
email & VARCHAR(254) & Да & Адрес электронной почты (уникален) \\ \hline
\dbfield{contact\_phone} &VARCHAR(20) & Нет & Контактный телефон \\ \hline
\dbfield{service\_status} &ENUM & Да & Статус обслуживания (active, limited, suspended, cancelled) \\ \hline
ip & VARCHAR(45) & Нет & IP-адрес абонента \\ \hline
\dbfield{installation\_address} &VARCHAR(255) & Нет & Адрес установки \\ \hline
\dbfield{gps\_coordinates} &VARCHAR(40) & Нет & Координаты установки \\ \hline
\dbfield{contract\_date} &DATE & Да & Дата заключения договора \\ \hline
password & VARCHAR(128) & Да & Хеш пароля абонента \\ \hline
\dbfield{deleted\_at} &DATETIME(6) & Нет & Отметка мягкого удаления
\end{xltabular}

Структура таблицы счетов \texttt{invoices} приведена в таблице~\ref{tbl:db_inv}.

\begin{xltabular}{\textwidth}{|p{3.0cm}|p{3.0cm}|p{1.6cm}|X|}
\caption{Структура таблицы \texttt{invoices}\label{tbl:db_inv}}\\ \hline
\centrow Поле & \centrow Тип & \centrow Обязат. & \centrow Назначение \\ \hline
\endfirsthead
\continuecaption{Продолжение таблицы \ref{tbl:db_inv}}
\centrow Поле & \centrow Тип & \centrow Обязат. & \centrow Назначение \\ \hline
\finishhead
id & BIGINT & Да & Первичный ключ \\ \hline
\dbfield{client\_id} &BIGINT & Да & Внешний ключ на \texttt{clients} \\ \hline
\dbfield{client\_plan\_id} &BIGINT & Нет & Внешний ключ на \texttt{clients\_plans} \\ \hline
\dbfield{invoice\_number} & \dbtype{VARCHAR(30)} & Да & Уникальный фискальный номер \\ \hline
\dbfield{issue\_date} &DATE & Да & Дата выпуска \\ \hline
\dbfield{due\_date} &DATE & Да & Срок оплаты \\ \hline
amount & \dbtype{DECIMAL(12,2)} & Да & Сумма без налога \\ \hline
\dbfield{tax\_amount} &DECIMAL(12,2) & Да & Сумма налога \\ \hline
\dbfield{total\_amount} &DECIMAL(12,2) & Да & Сумма с налогом \\ \hline
status & VARCHAR(20) & Да & Статус (draft, pending, paid, failed, cancelled) \\ \hline
\dbfield{payment\_method} &VARCHAR(20) & Нет & Метод оплаты \\ \hline
\dbfield{configuration\_snapshot} &JSON & Да & Снимок фискальных реквизитов на момент выпуска
\end{xltabular}

Структуры остальных прикладных таблиц (\texttt{plans}, \texttt{clients\_plans}, \texttt{wallets}, \texttt{transactions}, \texttt{tickets}, \texttt{messages}, \texttt{mikrotik\_routers}, \texttt{firewall\_filter\_rules}, \texttt{audits} и др.) построены по аналогичному принципу: суррогатный первичный ключ \texttt{id} типа \texttt{BIGINT}, внешние ключи на связанные сущности с заданным поведением при удалении (\texttt{CASCADE} / \texttt{SET~NULL} / \texttt{RESTRICT}), служебные поля \texttt{created\_at} и \texttt{updated\_at}, для критичных таблиц~--~поле \texttt{deleted\_at} для мягкого удаления. Состав полей каждой таблицы соответствует обобщенному описанию из таблицы~\ref{tbl:db_summary} и подробно проиллюстрирован структурами таблиц~\ref{tbl:db_clients} и~\ref{tbl:db_inv} как представительными примерами.

Назначение служебных таблиц фреймворка Laravel: \texttt{personal\_access\_tokens}~--~хранилище токенов Sanctum; \texttt{jobs}, \texttt{failed\_jobs}~--~очередь фоновых задач; \texttt{cache}, \texttt{sessions}~--~хранилища кэша и сеансов; \texttt{migrations}~--~журнал примененных миграций; \texttt{password\_reset\_tokens}~--~токены восстановления пароля. Для ускорения поиска объявлены индексы по \texttt{full\_name}, \texttt{document\_id} в \texttt{clients}; по \texttt{name} в \texttt{plans} и \texttt{mikrotik\_routers}; по \texttt{invoice\_number}, \texttt{client\_id}, \texttt{status} в \texttt{invoices}; по \texttt{slug} в \texttt{roles} и \texttt{permissions}. Первичные и внешние ключи индексируются автоматически средствами Eloquent~ORM.

\subsection{Подсистема интеграции с маршрутизаторами MikroTik}

\subsubsection{Сценарии синхронизации простых очередей}

Подсистема синхронизации простых очередей маршрутизатора решает задачу согласования двух автономных источников данных~--~локальной базы данных подписок абонентов и таблицы \texttt{/queue/simple/} на самом маршрутизаторе. Источником истины с точки зрения системы является локальная база; маршрутизатор после синхронизации должен содержать тот же набор очередей, что описан в базе. Реализация подсистемы инкапсулирована в сервисе \texttt{MikroTikQueueSyncService}, основным методом которого является \texttt{sync(MikrotikRouter \$router, bool \$cleanup = true)}, выполняющий полный цикл синхронизации.

Алгоритм синхронизации построен по принципу <<вычислить разность~--~применить порциями>>. На первом этапе сервис загружает из локальной базы данных все активные подписки абонентов (\texttt{ClientPlan} со статусом \texttt{active} и непустым полем \texttt{ip\_address}), выбирает связанные тарифные планы и абонентов для получения параметров скорости и IP-адреса. На втором этапе сервис подключается к маршрутизатору через \texttt{MikroTikService}, выполняет команду \texttt{/queue/simple/print} с фильтром по комментарию (комментарии содержат префикс <<isp\_gestor:>>, что отделяет управляемые системой очереди от других, существующих на маршрутизаторе по иным причинам), сохраняет результат в виде ассоциативного массива по идентификатору \texttt{.id} очереди.

На третьем этапе сервис вычисляет три подмножества операций: \emph{добавить} (подписки, представленные в локальной базе, но не имеющие соответствующих очередей на маршрутизаторе), \emph{удалить} (очереди, существующие на маршрутизаторе, но не имеющие соответствующих подписок в базе~--~такие очереди условно называются <<осиротевшими>>), \emph{обновить} (подписки и очереди с расходящимися параметрами скорости либо целевого IP-адреса).

На четвертом этапе сервис применяет вычисленные изменения к маршрутизатору. Применение выполняется порциями по 20 операций на одну пакетную команду, что снижает накладные расходы на инициализацию каждого отдельного запроса и при этом не превышает разумных пределов размера команды RouterOS. Для каждой операции фиксируется результат (успех либо текст ошибки) в журнале выполнения задачи; задача продолжает работу с оставшимися операциями даже при возникновении ошибок на отдельных шагах, что предотвращает зависание процесса синхронизации из-за единичного некорректного состояния.

Имя очереди формируется по шаблону \texttt{<<plan-\{plan\_id\}-cli-\{client\_id\}>>}; целевой объект (\texttt{target}) задается как IP-адрес абонента; максимальная скорость (\texttt{max-limit}) формируется в формате \texttt{<<\{up\}M/\{down\}M>>}, где \texttt{up} и \texttt{down}~--~скорости из тарифного плана в мегабитах в секунду; комментарий (\texttt{comment}) фиксируется как \texttt{<<isp\_gestor:active:plan=\{slug\}>>}, что в дальнейшем используется при чтении состояния маршрутизатора для отличия управляемых очередей от прочих.

Опциональный параметр \texttt{\$cleanup} управляет поведением подмножества <<удалить>>: при значении \texttt{true} осиротевшие очереди удаляются с маршрутизатора, что приводит к полному соответствию маршрутизатора локальной базе; при значении \texttt{false} удаление пропускается, и на маршрутизаторе могут остаться очереди, не имеющие соответствия в базе (например, очереди, созданные администратором вручную и не подлежащие управлению системой). Поведение по умолчанию~--~удаление осиротевших, поскольку именно в этом режиме обеспечивается полная синхронизация.

\subsubsection{Управление правилами брандмауэра}

Подсистема управления правилами брандмауэра отличается от подсистемы синхронизации очередей более сложной моделью изменения состояния: правила брандмауэра не отображаются на сущности абонентской базы автоматически (как очереди отображаются на подписки), поэтому источником истины для них является не локальная база, а решения сетевого администратора. Локальная база в этом случае выступает как <<черновик>>~--~рабочая копия конфигурации, в которую сетевой администратор вносит изменения через интерфейс панели администрирования, после чего пакет изменений атомарно применяется к маршрутизатору.

Подсистема реализована в сервисе \texttt{FirewallService} с четырьмя основными операциями: \texttt{loadFromRouter()}~--~загрузка текущего состояния правил с маршрутизатора в локальную базу; \texttt{validatePending()}~--~валидация локальных изменений до применения; \texttt{apply()}~--~атомарное применение локальных изменений к маршрутизатору; \texttt{rollback(\$applyLogId)}~--~откат состояния маршрутизатора к снимку из указанной записи журнала применений.

\emph{Операция \texttt{loadFromRouter}} выполняет команды \texttt{/ip/firewall/filter/print} и \texttt{/ip/firewall/nat/print}, преобразует результат в локальные модели \texttt{FirewallFilterRule} и \texttt{FirewallNatRule}, сохраняет идентификаторы правил на маршрутизаторе в поле \texttt{external\_id}. Каждое локальное правило при этом получает собственный UUID, который сохраняется неизменным при последующих изменениях и используется как стабильный идентификатор при отображении в интерфейсе.

\emph{Операция \texttt{validatePending}} обходит локальные правила со статусом <<черновик>> и проверяет совокупный набор изменений на согласованность: отсутствие дублирующих правил с одинаковыми условиями совпадения, корректность значений цепочек и действий, отсутствие ссылок на несуществующие интерфейсы, отсутствие циклических переходов между пользовательскими цепочками.

\emph{Операция \texttt{apply}} выполняется в трех фазах. Фаза~1~--~снимок <<до>>: сервис загружает с маршрутизатора актуальное состояние правил и сохраняет его в поле \texttt{snapshot\_before} новой записи \texttt{FirewallApplyLog}. Фаза~2~--~применение: сервис последовательно выполняет команды \texttt{/ip/firewall/filter/add}, \texttt{/ip/firewall/filter/set}, \texttt{/ip/firewall/filter/remove} (а также соответствующие команды цепочки \texttt{nat}) в порядке, обеспечивающем непрерывность сетевой связности (правила, разрешающие трафик, добавляются перед удалением правил, ограничивающих этот трафик). Фаза~3~--~снимок <<после>>: сервис загружает итоговое состояние и сохраняет его в поле \texttt{snapshot\_after} записи журнала.

При возникновении ошибки на фазе~2 (например, при потере связи с маршрутизатором либо при отказе маршрутизатора в применении правила) сервис автоматически выполняет аварийный откат: применяет снимок \texttt{snapshot\_before} к маршрутизатору, фиксирует ошибку в журнале выполнения задачи, возвращает пользователю ошибку через программный интерфейс. Такая схема обеспечивает атомарность применения: либо все изменения успешно применены, либо состояние маршрутизатора остается прежним.

\subsubsection{Откат изменений правил брандмауэра}

Операция отката \texttt{rollback(\$applyLogId)} принимает идентификатор записи журнала применений и применяет снимок \texttt{snapshot\_before} из этой записи к маршрутизатору. Откат сам по себе является применением, поэтому создает новую запись в журнале применений с пометкой \texttt{<<откат записи №...>>} и снимками \texttt{snapshot\_before} и \texttt{snapshot\_after}, отражающими состояние до и после отката. Эта рекурсивная природа журнала позволяет в дальнейшем откатить и сам откат, что приводит к восстановлению исходного состояния, действовавшего до отката.

В интерфейсе панели администрирования каждая запись журнала применений сопровождается кнопкой <<Откатить>>, доступной пользователям с правом \texttt{mikrotik.gestionar}. Нажатие на кнопку отображает модальное окно подтверждения с детальным сравнением снимков <<до>> и <<после>> применения; пользователь видит, какие именно правила будут изменены откатом, и подтверждает операцию. Подтверждение приводит к вызову операции \texttt{rollback} соответствующего сервиса.

\subsubsection{Управление подключением к маршрутизатору}

Сервис \texttt{MikroTikService} инкапсулирует низкоуровневое взаимодействие с маршрутизаторами и предоставляет вышестоящим сервисам и контроллерам типизированный интерфейс выполнения команд RouterOS~API. Конструктор сервиса принимает объект модели \texttt{MikrotikRouter} и выполняет следующие операции:
\begin{itemize}
\item расшифровывает пароль учетной записи маршрутизатора с применением facade \texttt{Crypt} Laravel и ключа \texttt{APP\_KEY};
\item конфигурирует объект клиента библиотеки \texttt{evilfreelancer/routeros-api-php} с параметрами хоста, порта, имени пользователя и расшифрованного пароля;
\item устанавливает таймаут подключения~--~5~секунд для первоначального соединения, 30~секунд для выполнения каждой отдельной команды;
\item включает поддержку TLS, если на стороне маршрутизатора указан порт 8729.
\end{itemize}

При завершении работы сервиса (автоматически при разрушении объекта) выполняется закрытие соединения через метод \texttt{disconnect()}. Это предотвращает накопление зависших TCP-соединений на стороне маршрутизатора, что особенно важно при выполнении большого числа коротких операций (например, при синхронизации нескольких сотен очередей подряд).

Команда \texttt{php artisan mikrotik:test \{router\_id\}} обеспечивает проверку доступности маршрутизатора и корректности учетных данных. Команда выполняет следующие операции в указанном порядке: подключение к маршрутизатору с зафиксированным таймаутом; выполнение команд \texttt{/system/identity/print} и \texttt{/system/resource/print} для получения базовой информации об устройстве; создание тестовой очереди \texttt{/queue/simple/add name=isp\_gestor\_test target=192.0.2.1 max-limit=1k/1k}; удаление тестовой очереди \texttt{/queue/simple/remove [find name="isp\_gestor\_test"]}; запись результата каждой операции в журнал выполнения. Команда применяется при первичной настройке нового маршрутизатора и при диагностике сбоев интеграции.

\subsection{Подсистема автоматизации}

\subsubsection{Очередь фоновых задач Laravel}

Часть операций системы (синхронизация маршрутизатора, формирование счетов, рассылка уведомлений, обработка платежных вебхуков) являются длительными по сравнению со временем отклика, ожидаемым от HTTP-запроса. Выполнение таких операций в рамках обработки запроса привело бы к блокировке клиентского интерфейса на десятки секунд, что недопустимо. Поэтому такие операции выносятся в \emph{очередь фоновых задач}: контроллер диспатчит соответствующую задачу в очередь и немедленно возвращает клиенту ответ; задача выполняется отдельным рабочим процессом (worker) асинхронно по отношению к породившему ее запросу.

Подсистема очередей Laravel поддерживает несколько драйверов: \texttt{sync} (синхронное выполнение в момент диспатча~--~для разработки и тестирования), \texttt{database} (хранение задач в таблице \texttt{jobs} базы данных), \texttt{redis} (хранение в Redis), \texttt{sqs} (Amazon Simple Queue Service), \texttt{beanstalkd} (Beanstalkd), \texttt{rabbitmq} (RabbitMQ через сторонний пакет). В разрабатываемой системе используется драйвер \texttt{database}, что упрощает развертывание (не требуется отдельный процесс Redis) при умеренном объеме задач (до нескольких тысяч в день).

Рабочий процесс запускается командой \texttt{php artisan queue:work}, которая в режиме демона периодически опрашивает таблицу \texttt{jobs} на наличие непринятых задач и выполняет их в порядке диспетчеризации с учетом приоритетов. Поддерживается одновременный запуск нескольких рабочих процессов для параллельного выполнения задач; в этом случае каждая задача захватывается ровно одним процессом через механизм оптимистических блокировок.

В системе определены четыре очереди с разными приоритетами:
\begin{itemize}
\item \texttt{high}~--~задачи, требующие немедленного выполнения (синхронизация очередей при приостановке абонента, отправка уведомлений о критических событиях);
\item \texttt{default}~--~задачи общего назначения (формирование счетов, обработка платежей);
\item \texttt{low}~--~задачи фонового характера, отложенные во времени (массовая рассылка, переиндексация данных для поиска);
\item \texttt{notifications}~--~отдельная очередь для рассылки уведомлений, изолированная от прочих, чтобы сбой провайдера уведомлений не блокировал основной поток задач.
\end{itemize}
Рабочий процесс запускается с указанием перечня обслуживаемых очередей в порядке убывания приоритета: \texttt{php artisan queue:work --queue=high,default,notifications,low}.

\subsubsection{Задача автоматического формирования счетов}

Задача \texttt{GenerateMonthlyInvoices} реализует прецедент <<Автоматическое формирование счетов по расписанию>> из подраздела~2.5. Метод \texttt{handle()} задачи выполняет следующие шаги:
\begin{enumerate}
\item загружает из базы все активные подписки абонентов (\texttt{ClientPlan::active()->get()});
\item для каждой подписки определяет дату следующего выпуска счета (поле \texttt{next\_billing\_date});
\item если \texttt{next\_billing\_date} наступила или прошла, проверяет наличие уже выпущенного счета за соответствующий период (поиск в таблице \texttt{invoices} по \texttt{client\_plan\_id} и диапазону дат выпуска);
\item при отсутствии счета формирует новый: вычисляет номер счета через сервис \texttt{InvoiceNumberingService} с учетом фискальной юрисдикции, фиксирует фискальный снимок через сервис \texttt{FiscalConfigSnapshotService}, создает запись \texttt{Invoice};
\item обновляет \texttt{next\_billing\_date} подписки на следующий расчетный период;
\item диспатчит задачу автоматического списания (\texttt{ChargeInvoice}) с идентификатором созданного счета.
\end{enumerate}

Идемпотентность задачи обеспечивается шагом 3: повторный запуск задачи в пределах одного расчетного периода не приводит к появлению дублирующих счетов, поскольку счет за этот период уже существует. Это особенно важно при сбоях рабочего процесса или планировщика, после которых задача может быть повторно поставлена в очередь.

\subsubsection{Задача автоматической приостановки услуг}

Задача \texttt{ProcessClientSuspension} запускается ежедневно по расписанию (по умолчанию в 08:00) и выполняет двухэтапную приостановку обслуживания абонентов с просроченной задолженностью. Этап~1~--~ограничение (перевод в статус \texttt{limited}): обрабатывает абонентов, у которых имеются неоплаченные счета со сроком оплаты, истекшим более \texttt{grace\_days} дней назад (значение из таблицы системных настроек, по умолчанию 5 дней). Этап~2~--~приостановка (перевод в статус \texttt{suspended}): обрабатывает абонентов, у которых аналогичная просрочка превышает \texttt{strict\_days} дней (по умолчанию 15 дней).

На каждом из этапов задача обрабатывает соответствующее множество абонентов следующим образом:
\begin{enumerate}
\item изменяет статус обслуживания абонента (\texttt{Client::service\_status});
\item создает запись события \texttt{ClientEvent} с типом \texttt{limited} или \texttt{suspended};
\item диспатчит фоновую задачу применения изменений на маршрутизаторе (\texttt{ApplyClientStatusChange}), которая удаляет очередь абонента и добавляет его IP-адрес в соответствующий адресный список (\texttt{limited} или \texttt{suspended});
\item опционально отправляет уведомление через подсистему уведомлений (электронное письмо, сообщение мессенджера, push-уведомление).
\end{enumerate}

Все операции этапа обрабатываются в одной транзакции базы данных, что гарантирует согласованность статуса абонента с записью события и со списком запланированных к выполнению задач.

\subsubsection{Задача автоматического восстановления услуг}

Задача \texttt{ProcessAutoReactivation} диспатчится не по расписанию, а по факту наступления значимого события~--~зачисления средств на лицевой счет абонента, находящегося в состоянии ограниченного или приостановленного обслуживания. Обработчик задачи выполняет следующие шаги:
\begin{enumerate}
\item загружает абонента и его текущие неоплаченные счета;
\item вычисляет суммарную задолженность как разницу между общей суммой неоплаченных счетов и текущим балансом лицевого счета;
\item если задолженности не осталось (баланс покрывает все просроченные счета), последовательно применяет оплату ко всем неоплаченным счетам в порядке нарастания даты выпуска;
\item изменяет статус обслуживания абонента на \texttt{active};
\item создает запись события \texttt{ClientEvent} с типом \texttt{reactivated};
\item диспатчит фоновую задачу пересоздания очереди MikroTik с параметрами актуального тарифа абонента.
\end{enumerate}

Такая модель автоматического восстановления обеспечивает мгновенное возобновление обслуживания при пополнении лицевого счета абонентом, что существенно улучшает пользовательский опыт по сравнению с моделью, в которой восстановление выполняется только по расписанию следующей проверки.

\subsubsection{Динамический планировщик задач}

Запуск задач по расписанию в Laravel выполняется планировщиком, описанным в файле \texttt{routes/console.php} (либо в методе \texttt{schedule()} класса \texttt{App\textbackslash Console\textbackslash Kernel}). Планировщик запускается ежеминутно через системный \texttt{cron} (запись \texttt{* * * * * cd /var/www/isp\_gestor \&\& php artisan schedule:run >> /dev/null 2>\&1}), на каждом запуске анализирует список объявленных задач и запускает те, для которых наступило расписание.

В разрабатываемой системе расписание задач не зашито жестко в код, а считывается из таблицы \texttt{automation\_settings} при каждом запуске планировщика через сервис \texttt{AutomationSettingsService}. Сервис возвращает массив активных задач с их расписаниями, который планировщик использует для регистрации соответствующих вызовов. Изменение расписания задачи через интерфейс панели администрирования вступает в силу при следующем запуске планировщика (то есть в течение одной минуты) без необходимости перезапуска фонового процесса.

Поддерживается специальный режим тестирования: при установке системной настройки \texttt{SCHEDULE\_TEST\_MODE = true} все задачи запускаются каждые 5 минут независимо от объявленного расписания, что используется при отладке для ускоренного срабатывания циклов биллинга. В производственной конфигурации этот режим должен быть отключен.

Дополнительно для разработки и тестирования предусмотрены консольные команды ручного запуска каждой из автоматических задач: \texttt{php artisan billing:generate-invoices}, \texttt{php artisan billing:process-suspensions}, \texttt{php artisan billing:check-reactivate}, \texttt{php artisan mikrotik:sync-queues [--cleanup]}. Команды позволяют запустить задачу немедленно (без ожидания расписания) с теми же параметрами, что и при автоматическом срабатывании, и проверить корректность ее работы на тестовых данных.

\subsection{Контрольный пример работы системы}

В настоящем подразделе приведены экранные формы основных рабочих мест системы, иллюстрирующие реализацию ключевых прецедентов из подраздела~2.5. Изображения подготовлены на основе работающего экземпляра системы, развернутого в стендовых условиях с подключенным маршрутизатором MikroTik RB951Ui-2HnD.

На рисунке~\ref{fig:scr_dashboard} представлена главная страница панели администрирования с ключевыми показателями деятельности оператора.

\begin{figure}[H]
\centering
\includegraphics[width=\linewidth,keepaspectratio]{Dashboard}
\caption{Главная страница панели администрирования}
\label{fig:scr_dashboard}
\end{figure}

На рисунке~\ref{fig:scr_clients} представлен реестр абонентов.

\begin{figure}[H]
\centering
\includegraphics[width=\linewidth,keepaspectratio]{Clientes}
\caption{Реестр абонентов с детальной карточкой}
\label{fig:scr_clients}
\end{figure}

На рисунке~\ref{fig:scr_firewall} представлена страница управления правилами брандмауэра MikroTik.

\begin{figure}[H]
\centering
\includegraphics[width=\linewidth,keepaspectratio]{CortaFuegos}
\caption{Управление правилами брандмауэра MikroTik}
\label{fig:scr_firewall}
\end{figure}

На рисунке~\ref{fig:scr_chat} представлена подсистема технической поддержки.

\begin{figure}[H]
\centering
\includegraphics[width=\linewidth,keepaspectratio]{SoporteTecnico}
\caption{Подсистема технической поддержки с внутренней перепиской сотрудников}
\label{fig:scr_chat}
\end{figure}

Полный комплект экранных форм всех разделов панели администрирования вынесен в приложение к настоящей пояснительной записке.
